% T23 source: simple_litreview.tex lines 25--41
\subsection{Scheduling Theory and Repeated Shared Service}

The mathematical core of levels M1--M2 is standard. For identical parallel
machines, the conventional quantities are makespan, total workload, the
duration of the longest job, and the critical path
\cite{pinedo2008scheduling}; list scheduling under precedence constraints and
the LPT rule for independent jobs have classical guarantees
\cite{graham1966bounds,graham1969bounds}. Disjunctive arcs, or resource-ordering
arcs, are likewise a standard way to represent competition among machines
\cite{pinedo2008scheduling}. The article therefore does not claim novelty for
makespan, a directed acyclic graph (DAG), the critical path, a
resource-augmented graph, or a shared resource.

The resource-constrained project scheduling problem (RCPSP) combines a DAG of
precedence relations with renewable resources of finite capacity, while
multi-mode formulations permit alternative duration and resource profiles
\cite{brucker1999resource}. Within such a core, attention can be represented as
a resource of capacity 1 and phases as separate operations. The standard model,
however, does not specify a coefficient $k$ for a software task and operating
mode, the meaning of an accepted result, an operational decomposition
$a_i/h_i$, or internal integration and attention-switching overhead.

Still closer analogues are parallel-machine models with a common server. In
setup-only formulations, the server sequentially prepares jobs before their
parallel processing \cite{hall2000parallel,kravchenko1997parallel}. In a
loading/unloading formulation, each job uses the same server before and after
machine processing; a requirement for immediate unloading may keep a machine
occupied after processing until the server becomes available
\cite{elidrissi2023loading}. This is a formal analogue of the sequence ``human
task specification---agent work---human review'' and a partial analogue of
slot retention, but it lacks a DAG and a repeated ``agent rework---human
review'' cycle.

An even closer reentrant-scheduling formulation routes a job through its
assigned primary machine, a remote server, and then back to the same primary
machine \cite{chakhlevitch2009reentrant}. It reproduces the alternation
``worker---shared service---the same worker,'' but includes only one repeated
visit to the server, independent jobs, and a prespecified route. Importantly,
the primary machine is released during remote service; this model therefore
does not reproduce M4 local blocking. Neither common-server models nor the
reentrant formulation by themselves include a coefficient $k$ for software
work, human acceptance, or integration overhead.

This comparison clarifies the contribution of M4: established scheduling
mechanisms are used as the core, while the proposed mapping to the software
process links the complete cycle ``human task specification---agent
work---human review---agent rework---human acceptance'' to the task's operating
mode, separate phase durations, retention of the assigned agent slot, and the
fact that the result has been accepted. If a human-attention phase does not
retain the slot, $A$ and $H$ should be scheduled as separate resources; this
corresponds to a different resource-use rule.

\subsection{Human Attention in Supervising Autonomous Workers}

Research on human--robot interaction (HRI) provides a substantive structural
analogue for one operator and multiple autonomous workers. The fan-out measure,
the number of robots per operator, is related to the durations of interaction
and autonomous work without operator attention \cite{olsen2004fanout}; models
of multitask control require service for one robot to fit within the autonomous
intervals of the others \cite{crandall2005validating}. State-aware attention
allocation shows that the outcome depends on the selection rule and system
state, not only on the average fraction of involvement
\cite{crandall2011computing}. Threaded-cognition theory further cautions that
not all human activity is uniformly sequential: interference arises at
specific cognitive resources \cite{salvucci2008threaded}.

Cummings and Mitchell tested an extended robots-per-operator model in a
simulation experiment with 12 participants and multiple unmanned aerial
vehicles (UAVs) \cite{cummings2008predicting}. Accounting for waiting for
interaction, waiting in the queue, waiting to regain situational awareness, and
cognitive reorientation reduced the model-based estimate of operator throughput
by 36--67\% under the conditions examined; the authors explicitly emphasize
that the result depends on the task and context. The experiment supports a
sequential service queue and a finite useful number of autonomous workers only
as a structural analogy. It does not calibrate the number of coding agents,
$\gamma(P)$, or the makespan of software work: the vehicles were homogeneous
and independent, and the outcome was a model-based estimate of throughput, not
the completion time of an accepted software task.
